% Cover letter for the submission of NewMethod.tex to the
% Journal of Computational Algebra (Elsevier).  Submitted as a separate
% file through Editorial Manager (editorialmanager.com/jaca).
%
% TO FILL IN BEFORE SENDING:
%   - check the guest editors' titles in the salutation (all three are
%     professors, but confirm before sending);
%   - the archival deposit's DOI in the "Code and data" paragraph, once
%     the Zenodo (or equivalent) record exists;
%   - the date.

\documentclass[11pt]{article}

\usepackage[margin=0.9in]{geometry}
\usepackage[T1]{fontenc}
\usepackage{url}
\usepackage[hidelinks]{hyperref}

\setlength{\parindent}{0pt}
\setlength{\parskip}{4pt}
\pagestyle{empty}

\begin{document}

% The letter runs a few lines past a page; give page 1 the extra depth
% rather than spill a signature block onto a second page.
\enlargethispage{5\baselineskip}

\hfill Brent Baccala \quad \url{cosine@freesoft.org}

\hfill \today

\vspace{6pt}

Dear Professors Guo, Chen and Ovchinnikov,

I am submitting {\it A New Method of Solving PDEs} to the {\it Journal of
Computational Algebra}, for the special issue {\it Differential Algebra and
Related Topics}, as a Full-Length Research Article.

{\bf What the paper does.}  It gives an algorithmic method for finding exact
solutions of partial differential equations, and of systems of them, by fitting
the system to a parameterized {\it ansatz}: a finite set of differential
polynomials whose coefficients are undetermined constants.  The method computes,
inside the space of those constants, the locus at which the ansatz solves the
system.  The computation is a differential Thomas decomposition of the ansatz,
Ritt's full reduction of each equation of the target system modulo each
component, and a projection of the remainders onto the constants, assembled with
constructible-set operations.  A completeness theorem establishes that what the
algorithm returns is exactly the locus sought.

{\bf Fit with the call.}  The paper is differential algebra in the broader
sense the call describes --- the algebraic theory of differential equations,
applied to a problem in the sciences --- and it is squarely {\it about}
computation, which the call particularly solicits: the theory exists to justify
two algorithms, and the algorithms were run.  It also meets the journal's
general criterion, that both the mathematics and the algorithms contribute to
the result.  Section~3 develops the theory: specialization of a
simple system at a constant vector, the three loci and their containments, and
the completeness theorem.  The two algorithms are
given as pseudocode with proofs of what they return.  Section~5 is a
computation, carried out with an implementation, that produced a previously
unknown closed-form solution of the hydrogen Schr\"odinger equation,
$\Psi = J_0\bigl(2\sqrt{x+r}\bigr)$ at $E=0$, which is neither spherically nor
Cartesian separable.  It is not reachable by standard use of an existing
implementation, and it is not heuristic: the completeness theorem certifies that
nothing valid was missed.

{\bf What I claim is new.}  Three things.

The first is the embedding of ODEs in a {\it partial} differential ring with
arbitrary flexibility in the choice of the ODE's independent variable.  The
ansatz does not fix what the ODE is an equation in: the independent variable is
itself a parameterized expression --- in the hydrogen example,
$v = v_1 x + v_2 y + v_3 z + v_4 r$ --- and the constants that determine it are
solved for alongside the ODE's own coefficients.  This is what lets the method
reach solutions separable in no coordinate system chosen beforehand, the
parabolic-coordinate solution above among them.

The second is the consistency and membership loci themselves --- more
generally, any notion of a locus in the space of constants attached to a
parameterized ansatz.  I am not aware of prior work that defines such loci,
distinguishes them, or establishes the containments between them.  The paper defines
three, shows the containments are strict, and shows what each one answers.

The third is the pair of algorithms that compute the two loci, together with
the proofs that they do.

{\bf Declarations.}  I am the sole author.  The work has not been published
previously and is not under consideration elsewhere.  I have no competing
interests to declare, and the research received no specific grant from any
funding agency.

{\bf Generative AI.}  I used AI tools substantially in preparing this work, well
beyond checking language and readability, and would rather say so at the outset.
The manuscript carries the declaration section your guide requires, immediately
before the reference list, itemizing that use; a commit-level record is public in
the paper's repository at \url{https://github.com/BrentBaccala/NewMethod}.  All
of it was reviewed and edited by me, and I take full responsibility for it.

{\bf Code and data.}  The Sage scripts and the two supporting packages that
produced the paper's computations are described in Section~4 and, per the
journal's research-data policy, deposited in an archival repository and cited in
the reference list.  Thank you for considering the manuscript.

% Kept in a minipage so the sign-off cannot be split from the name.
\begin{minipage}{\linewidth}
Yours sincerely,

\vspace{18pt}

Brent Baccala
\end{minipage}

\end{document}
